\section*{Limitations}
This paper focuses exclusively on Wikipedia, the largest open text corpus. As a result, we do not explore other potentially valuable applications of corpus-level inconsistency detection, such as technical texts (e.g., academic, medical, or legal documents) or structured data sources like databases and knowledge graphs. We also leave detecting cross-lingual inconsistencies across different language versions of Wikipedia to future work.

\section*{Ethical Considerations}
We do not anticipate risks or ethical concerns arising from the publication of \dataset.

For crowdsourcing, we compensated annotators per task, with an overall rate of at least \$16 per hour. Participants in our user study were compensated at \$20 per hour. The study was approved by our institution's IRB, and participants provided informed consent. No personally identifiable information was collected during annotation or the user study. We release \dataset under a license compatible with Wikipedia's license.

Regarding computational resources, we used a CPU-based machine to serve the Wikipedia index and relied on commercial LLM APIs, making direct estimation of carbon footprint difficult. As a proxy, the total experimental cost did not exceed \$4,000.

\section*{Acknowledgment}
This work is supported in part by the Verdant Foundation, the Alfred P. Sloan Foundation, Microsoft Azure AI credits, and the NAIRR Pilot program.
